\frontmatter
\begin{titlepage}
\centering
\vspace*{30mm}
{\sffamily\bfseries\color{accent}\fontsize{40}{44}\selectfont Italian Learning Lab\par}
\vspace{10mm}
{\sffamily\LARGE Grammar, Reader and Workbook\par}
\vspace{6mm}
{\large From the first sounds to the complete verb system\par}
\vspace{4mm}
{\large\color{muted} Levels A0 to about B1--B2\par}
\vfill
{\large Companion to the Italian Learning Lab website\par}
\vspace{4mm}
{\color{muted}Edition generated from course data version @@version@@\par}
\vspace{20mm}
\end{titlepage}

\thispagestyle{empty}
\vspace*{\fill}
{\small
\textbf{Italian Learning Lab: Grammar, Reader and Workbook}

Written by the Italian Learning Lab contributors. The educational content (lessons, explanations, examples, exercises, vocabulary, stories and rhymes) is released under the Creative Commons Attribution 4.0 International licence (CC BY 4.0): you may copy, adapt and share it, including commercially, if you give appropriate credit. The software that generates this book is released under the MIT licence.

This book is generated automatically from the same data and the same grammar engine as the website, so the two always contain the same lessons, rules, exercises, answers, vocabulary and texts. Verb tables were produced by the conjugation engine and checked by its automated tests.

\textbf{Verification.} No native speaker, Italian teacher or linguist has reviewed this course. Statements about the passato remoto and its regional use, past-participle agreement, the auxiliary of modal verbs, the periodo ipotetico and the choice between congiuntivo and indicativo were compared with the reference works listed in Appendix~\ref{sources}; everything else is the authors' own work based on standard grammar. Appendix~\ref{sources} records, topic by topic, what has and has not been checked. Corrections are welcome.

Levels (A0, A1, A2, B1, B2) are approximate and describe intended difficulty only. This book is not affiliated with, or certified by, any examination body.
}
\clearpage

\chapter{How to use this book}

This book contains everything in the Italian Learning Lab website except listening and speaking: there is no audio, no read-aloud and no pronunciation practice. It can be used on its own, as a printed textbook and workbook, or next to the website. Where the website is interactive, the book prints the result instead: every story sentence comes with its translation and a verb analysis, every verb in the course is conjugated in full or listed with its principal parts, and every exercise has an answer and an explanation in the answer key.

\section*{The teaching principle}

The course teaches accurate, standard Italian from the first lesson. Difficulty is controlled by the choice of vocabulary, the length of sentences and the reading material --- not by leaving out grammatical distinctions or by stating rules that later have to be unlearned. Exceptions, alternative forms, differences of register and regional differences are stated where they occur. Every rule box says what kind of statement it is:

\begin{notebox}{rulecol}\notelabel{rulecol}{Rule}{}A grammatical rule of standard Italian. Departing from it is an error in the standard language.\end{notebox}
\begin{notebox}{tendcol}\notelabel{tendcol}{Tendency}{}A strong regularity, a stylistic preference or a difference of register or region. Alternatives exist and are not errors.\end{notebox}
\begin{notebox}{exccol}\notelabel{exccol}{Exception}{}A form or word that does not follow the rule just stated.\end{notebox}
\begin{notebox}{tipcol}\notelabel{tipcol}{Tip}{}A practical way to remember or apply something.\end{notebox}

\section*{What is in the book}

\begin{description}[leftmargin=0pt,itemsep=3pt]
\item[Part I, Course.] @@modules@@ chapters (one per module) with @@lessons@@ lessons, from spelling and pronunciation rules to the congiuntivo, the sequence of tenses, indirect speech and the passive. Each lesson has objectives, explanations, verb tables, worked examples, exercises, a summary and its related vocabulary.
\item[Part II, Reader.] @@stories@@ graded stories and @@rhymes@@ rhymes. Each story is printed first as continuous text, then sentence by sentence with a translation, notes and an analysis of every verb.
\item[Part III, Practice.] A diagnostic test to find your starting point and mixed revision sets.
\item[Part IV, Grammar reference.] An overview of the whole verb system, the sequence of tenses, the types of hypothetical sentence, the articulated prepositions, full conjugations of the auxiliaries, the regular models and every irregular verb in the course, and the principal parts of all @@verbs@@ verbs.
\item[Appendices.] The answer key for all @@exercises@@ exercises; the @@words@@ words of the course grouped by topic; sources and the verification register; an alphabetical index of Italian words.
\end{description}

\section*{Suggested ways of working}

\begin{itemize}
\item \textbf{In order.} Work through Part I chapter by chapter. After each lesson, do its exercises in writing, check them in the answer key, and read the story or rhyme suggested under \emph{Apply it}. Every lesson can be revisited at any time; nothing depends on finishing the previous one, but later lessons assume the earlier ones.
\item \textbf{From a placement.} Take the diagnostic test in Chapter~\ref{diag}. Start with the first chapter in which you made a mistake.
\item \textbf{As a reference.} Use Part IV and the index. Every tense in the overview table points to the lesson that explains it.
\item \textbf{As a reader.} The stories in Part II can be read before the grammar they use: the verb analysis explains every form and gives the section where it is taught.
\end{itemize}

\section*{Conventions}

Italian words in English text are in \textit{italics}. In verb tables the persons are \textit{io, tu, lui/lei, noi, voi, loro}; the formal \textit{Lei} and \textit{Loro} use the third-person forms. Endings are printed in \hlending{green} where the lesson asks you to notice them; forms that differ from the regular pattern of their family are printed in \irr{red-brown}. Where two forms are separated by a slash (\textit{credei/credetti}), both are standard. A cross-reference such as \S\ref{les:passato-remoto} points to a section of the book.

\section*{Exercises and answers}

Exercises are designed to be written out. Blanks are shown as lines; in matching exercises the right-hand column has been put in a different order; in word-order exercises the words are given in alphabetical order. The answer key gives every answer that the website accepts, the explanation, and, where the website anticipates a typical mistake, what is wrong with it. When checking your answers, accents, articles, endings and double consonants count; capital letters and final punctuation do not.

\section*{Reviewing vocabulary on paper}

The website schedules vocabulary with a Leitner system, which works equally well with paper cards. Write each word on a card (Italian with its article on one side, English on the other) and keep six boxes, reviewed after 1, 2, 4, 7, 14 and 30 days. A new card starts in a \emph{learning} pile that is reviewed every day. After each review: if you did not remember the word, the card goes back to the learning pile; if you remembered it with difficulty, it stays in its box and comes back after the interval of the previous box; if you remembered it, it moves up one box; if it was effortless, it moves up two boxes. A card in the last box stays there and comes back after 30 days. A word in box 4 or higher (an interval of a week or more) can be counted as known.

\tableofcontents
